\section{Electrical simulation}

\subsection{Scope and reference configuration}
This section documents the control cycle of the Organic Rankine Cycle (ORC)
power block and the electrical model derived from it for grid interconnection
studies. The reference configuration is a \textbf{40\,MW-class ORC unit,
directly grid-connected through a synchronous generator} --- the arrangement
used by large ORC plants (Ormat and Turboden class, roughly 0.3--20\,MW per
unit and above), in which the expander drives a conventional four-pole
synchronous machine at 1{,}500/1{,}800\,rpm through a reduction gearbox.

Small ORC units below roughly 300\,kW, and high-speed direct-drive machines,
instead use a permanent-magnet generator with a rectifier and a grid-tie
inverter; such a plant is modelled with the WECC generic inverter-based
resource models (\texttt{REGC\_A} / \texttt{REEC\_A} / \texttt{REPC\_A}) and is
outside the scope of this section.

The target interconnection is the \textbf{Western Interconnection (WECC)}, which
drives the model-acceptance and power system stabilizer requirements in
Section~\ref{sec:pss}.

\subsection{Plant architecture and control cycle}
Figure~\ref{fig:orc-control-cycle} shows the complete plant: the energy path
from heat source to point of interconnection, the closed working-fluid loop,
and the control and measurement layer wrapped around them.

\begin{figure}[htbp]
  \centering
  \resizebox{\linewidth}{!}{%
  \begin{tikzpicture}[
    font=\scriptsize,
    phys/.style={draw=atFg, line width=0.5pt, fill=atBgSubtle, rounded corners=1pt,
                 align=center, inner sep=2.5pt},
    ctrl/.style={draw=atAccentFg, line width=0.9pt, fill=white, rounded corners=1pt,
                 align=center, inner sep=2.5pt},
    dcsb/.style={draw=atAccentFg, line width=1.2pt, fill=atAccent!22, rounded corners=1pt,
                 align=center, inner sep=3pt},
    wf/.style={draw=atInfo, line width=0.9pt, -{Latex[length=1.8mm]}},
    wfp/.style={draw=atInfo, line width=0.9pt},
    brine/.style={draw=atDanger, line width=0.9pt, -{Latex[length=1.8mm]}},
    elec/.style={draw=atFg, line width=1.4pt, -{Latex[length=2mm]}},
    elecp/.style={draw=atFg, line width=1.4pt},
    shaft/.style={draw=atFgMuted, line width=1.8pt},
    sig/.style={draw=atAccentFg, dashed, line width=0.55pt, -{Latex[length=1.6mm]}},
    sigb/.style={draw=atAccentFg, dashed, line width=0.55pt,
                 {Latex[length=1.6mm]}-{Latex[length=1.6mm]}},
    lab/.style={font=\tiny, text=atFgMuted},
    sg/.style={font=\tiny, text=atAccentFg},
    chip/.style={circle, draw=atFg, fill=white, line width=0.4pt,
                 inner sep=0.7pt, font=\tiny},
    tap/.style={circle, fill=atAccentFg, inner sep=1.1pt},
  ]

  % ---------- physical blocks : energy path ----------
  \node[phys, minimum width=15mm, minimum height=13mm] (heat)  at (0,0)
        {HEAT\\SOURCE};
  \node[lab] at (0,-0.95) {brine / waste heat};
  \node[phys, minimum width=19mm, minimum height=13mm] (evap)  at (2.7,0)
        {PREHEATER\\+ EVAPORATOR};
  \node[phys, minimum width=9mm,  minimum height=9mm]  (gear)  at (7.9,0)
        {GEAR\\BOX};
  \node[phys, circle, minimum size=14mm] (gen) at (9.3,0) {G};
  \node[lab] at (9.3,-1.0) {synchronous generator};
  \node[phys, circle, minimum size=9mm] (t1) at (11.2,0) {};
  \node[phys, circle, minimum size=9mm] (t2) at (11.75,0) {};
  \node[lab] at (11.5,-0.85) {GSU XFMR};
  \node[phys, minimum width=5mm, minimum height=5mm] (brk) at (13.0,0) {};
  \node[lab] at (13.0,-0.6) {52G};
  \node[phys, minimum width=15mm, minimum height=10mm] (grid) at (15.0,0)
        {POI /\\GRID};

  % turbine (hand-drawn trapezium)
  \draw[draw=atFg, line width=0.5pt, fill=atBgSubtle]
        (5.55,0.42) -- (7.05,0.85) -- (7.05,-0.85) -- (5.55,-0.42) -- cycle;
  \node[align=center] at (6.30,0) {ORC\\TURBINE};

  % ---------- physical blocks : heat rejection loop ----------
  \node[phys, minimum width=34mm, minimum height=11mm] (cond) at (7.2,-2.4)
        {AIR-COOLED CONDENSER};
  \node[phys, circle, minimum size=12mm] (pump) at (4.0,-2.4) {FEED\\PUMP};
  \node[ctrl, minimum width=18mm, minimum height=8mm]  (fvfd) at (10.4,-1.9)
        {FAN VFD};
  \node[ctrl, minimum width=22mm, minimum height=6mm]  (pvfd) at (4.0,-3.9)
        {PUMP VFD};
  % fans
  \draw[draw=atFg, line width=0.5pt, fill=atBgSubtle] (7.4,-1.68) circle (0.18);
  \draw[draw=atFg, line width=0.4pt] (7.27,-1.81) -- (7.53,-1.55);
  \draw[draw=atFg, line width=0.4pt] (7.53,-1.81) -- (7.27,-1.55);
  \draw[draw=atFg, line width=0.5pt, fill=atBgSubtle] (8.6,-1.68) circle (0.18);
  \draw[draw=atFg, line width=0.4pt] (8.47,-1.81) -- (8.73,-1.55);
  \draw[draw=atFg, line width=0.4pt] (8.73,-1.81) -- (8.47,-1.55);

  % ---------- control blocks ----------
  \node[dcsb, minimum width=46mm, minimum height=9mm] (dcs) at (9.0,5.0)
        {PLANT DCS / PLANT CONTROLLER\\[1pt]
         \tiny dispatch $\cdot$ sequencing $\cdot$ POI schedule};
  \node[ctrl, minimum width=20mm, minimum height=8mm] (gov)  at (4.3,3.0)
        {GOVERNOR\\[1pt]\tiny speed / load};
  \node[ctrl, minimum width=20mm, minimum height=8mm] (avr)  at (7.2,3.0)
        {AVR + PSS\\[1pt]\tiny excitation};
  \node[ctrl, minimum width=18mm, minimum height=8mm] (sync) at (11.6,3.0)
        {AUTO-SYNC\\[1pt]\tiny 25 sync-check};
  \node[ctrl, minimum width=22mm, minimum height=8mm] (prot) at (14.3,3.0)
        {PROTECTION\\[1pt]\tiny 87G$\cdot$32$\cdot$40$\cdot$81$\cdot$27/59};

  % ---------- brine loop ----------
  \draw[brine] ([yshift=2.5mm]heat.east)  -- ([yshift=2.5mm]evap.west);
  \draw[brine] ([yshift=-2.5mm]evap.west) -- ([yshift=-2.5mm]heat.east);
  \node[lab] at (2.05,0.42) {hot};
  \node[lab] at (2.05,-0.42) {reinject};

  % ---------- working fluid loop ----------
  \draw[wf]  (evap.east) -- (5.55,0);
  \draw[fill=white, draw=atInfo, line width=0.7pt]
        (4.15,0.22) -- (4.15,-0.22) -- (4.55,0.22) -- (4.55,-0.22) -- cycle;
  \draw[fill=white, draw=atInfo, line width=0.7pt]
        (4.95,0.22) -- (4.95,-0.22) -- (5.35,0.22) -- (5.35,-0.22) -- cycle;
  \node[lab] at (4.35,-0.45) {ESV};
  \node[lab] at (5.15,0.45)  {TCV};
  \draw[wf]  (6.30,-0.64) -- (6.30,-1.85);
  \node[lab, anchor=west] at (6.40,-1.25) {low-p vapour};
  \draw[wf]  (cond.west) -- (4.62,-2.4);
  \draw[wf]  (3.38,-2.4) -- (2.7,-2.4) -- (2.7,-0.66);
  \node[lab, anchor=west] at (2.80,-1.6) {liquid};

  % ---------- shaft ----------
  \draw[shaft] (7.05,0) -- (gear.west);
  \draw[shaft] (gear.east) -- (gen.west);

  % ---------- electrical ----------
  \draw[elecp] (gen.east) -- (t1.west);
  \draw[elecp] (t2.east)  -- (brk.west);
  \draw[elec]  (brk.east) -- (grid.west);

  % ---------- measurement taps ----------
  \node[tap] (spd)  at (8.47,0)  {};
  \node[tap] (vtct) at (10.15,0) {};
  \node[tap] (ctp)  at (12.5,0)  {};
  \node[tap] (poi)  at (13.8,0)  {};
  \node[tap] (pev)  at (2.2,-0.65) {};
  \node[tap] (pcd)  at (8.90,-2.4) {};

  % ---------- control signals ----------
  % DCS -> governor
  \draw[sig] (7.0,4.55) -- (7.0,4.25) -- (4.3,4.25) -- (4.3,3.4);
  \node[sg, anchor=west] at (4.45,4.40) {MW setpoint};
  % DCS -> AVR
  \draw[sig] (8.2,4.55) -- (8.2,3.75) -- (7.6,3.75) -- (7.6,3.4);
  \node[sg, anchor=west] at (8.30,3.95) {V / Q setpoint};
  % POI -> DCS
  \draw[sig] (13.8,-0.12) -- (13.8,-0.95) -- (16.2,-0.95) -- (16.2,5.95)
             -- (9.8,5.95) -- (9.8,5.45);
  \node[sg, anchor=east] at (16.0,6.15) {MW $\cdot$ MVAr $\cdot$ V at POI};
  % governor -> TCV
  \draw[sig] (3.7,2.6) -- (3.7,1.2) -- (5.15,1.2) -- (5.15,0.26);
  \node[sg, anchor=west] at (3.80,1.35) {valve command};
  % speed -> governor
  \draw[sig] (8.47,0.12) -- (8.47,2.0) -- (4.7,2.0) -- (4.7,2.6);
  \node[sg, anchor=west] at (5.60,2.15) {shaft speed $n$};
  % AVR -> field
  \draw[sig] (7.2,2.6) -- (7.2,2.32) -- (9.3,2.32) -- (9.3,0.72);
  \node[sg, anchor=west] at (7.90,2.47) {field current};
  % VT/CT -> AVR
  \draw[sig] (10.15,0.12) -- (10.15,3.0) -- (avr.east);
  \node[sg, anchor=south] at (9.30,3.05) {VT / CT};
  % sync <-> breaker
  \draw[sigb] (11.6,2.6) -- (11.6,1.0) -- (13.0,1.0) -- (13.0,0.3);
  \node[sg, anchor=west] at (11.70,1.15) {V, $f$, $\varphi$ / close};
  % sync -> governor + AVR bias
  \draw[sig] (11.0,3.4) -- (11.0,4.05) -- (4.9,4.05) -- (4.9,3.4);
  \node[sg, anchor=west] at (6.30,3.90) {speed / voltage match bias};
  % CT -> protection
  \draw[sig] (12.5,0.12) -- (12.5,1.6) -- (14.3,1.6) -- (14.3,2.6);
  % protection -> trip breaker
  \draw[sig] (14.8,2.6) -- (14.8,0.55) -- (13.15,0.55) -- (13.15,0.25);
  \node[sg, anchor=west] at (14.90,0.75) {trip 52G};
  % protection -> ESV
  \draw[sig] (13.8,2.6) -- (13.8,1.72) -- (4.35,1.72) -- (4.35,0.26);
  \node[sg, anchor=west] at (6.60,1.87) {emergency trip $\rightarrow$ ESV quick-close};
  % evaporator pressure -> pump VFD
  \draw[sig] (2.2,-0.77) -- (2.2,-3.9) -- (2.9,-3.9);
  \node[sg, anchor=east] at (2.10,-2.9) {$p_{\mathrm{evap}}$};
  \draw[sig] (4.0,-3.6) -- (4.0,-3.02);
  % condenser pressure -> fan VFD
  \draw[sig] (9.02,-2.4) -- (10.4,-2.4) -- (10.4,-2.3);
  \node[sg, anchor=north] at (9.70,-2.48) {$p_{\mathrm{cond}}$};
  \draw[sig] (9.5,-1.68) -- (8.82,-1.68);

  % ---------- item numbers (keyed to Table 1) ----------
  \node[chip] at (-0.75,0.65) {1};
  \node[chip] at (1.75,0.65)  {2};
  \node[chip] at (4.35,0.45)  {3};
  \node[chip] at (5.60,0.60)  {4};
  \node[chip] at (7.45,0.45)  {5};
  \node[chip] at (8.70,0.60)  {6};
  \node[chip] at (10.80,0.42) {7};
  \node[chip] at (14.30,0.50) {8};
  \node[chip] at (5.55,-1.85) {9};
  \node[chip] at (3.42,-1.92) {10};

  % ---------- legend ----------
  \draw[wfp]   (-0.9,-5.1) -- (0.1,-5.1);
  \node[lab, anchor=west] at (0.2,-5.1) {working fluid};
  \draw[draw=atDanger, line width=0.9pt] (2.8,-5.1) -- (3.8,-5.1);
  \node[lab, anchor=west] at (3.9,-5.1) {heat-source brine};
  \draw[shaft] (6.7,-5.1) -- (7.7,-5.1);
  \node[lab, anchor=west] at (7.8,-5.1) {shaft (mechanical)};
  \draw[elecp] (10.7,-5.1) -- (11.7,-5.1);
  \node[lab, anchor=west] at (11.8,-5.1) {electrical power};
  \draw[draw=atAccentFg, dashed, line width=0.55pt] (14.2,-5.1) -- (15.2,-5.1);
  \node[lab, anchor=west] at (15.3,-5.1) {control / measurement};
  \end{tikzpicture}}
  \caption{ORC plant control cycle. Energy flows left to right; the working
  fluid circulates evaporator $\rightarrow$ turbine $\rightarrow$ condenser
  $\rightarrow$ feed pump. Numbers key to Table~\ref{tab:orc-parts}. Setpoint
  links from the plant DCS to the pump and fan VFDs, and the recuperator common
  on many ORC cycles, are omitted for legibility; neither adds a control loop.}
  \label{fig:orc-control-cycle}
\end{figure}

\subsubsection{Physical parts of the energy path}

\begin{table}[htbp]
  \centering
  \renewcommand{\arraystretch}{1.25}
  \begin{tabularx}{\linewidth}{@{}c l X@{}}
    \toprule
    \textbf{\#} & \textbf{Part} & \textbf{Function and relevance} \\
    \midrule
    1  & Heat source loop      & Geothermal brine or waste-heat stream with its
                                 circulation pumps. Sets the total heat
                                 available and is the effective fuel valve of
                                 the plant. \\
    2  & Preheater, evaporator & Heats and vaporises the organic working fluid.
                                 Key measurements: evaporator pressure and
                                 superheat. \\
    3  & ESV and TCV           & Spring-closed emergency stop valve (tripped by
                                 protection) and the throttle valve modulated by
                                 the governor. The only fast handles on turbine
                                 power. \\
    4  & ORC turbine           & Axial or radial-inflow expander converting
                                 vapour enthalpy to shaft torque. Fitted with
                                 speed pickups and an independent overspeed
                                 trip. \\
    5  & Gearbox               & Reduces turbine speed to synchronous speed
                                 (1{,}500/1{,}800\,rpm). Omitted on machines
                                 designed to drive the generator directly. \\
    6  & Synchronous generator & Four-pole machine with shaft-mounted brushless
                                 exciter. Field current from the AVR sets
                                 terminal voltage and reactive power. Modelled
                                 as \texttt{GENROU}. \\
    7  & GSU transformer       & Steps generator voltage (typically
                                 11--13.8\,kV) up to interconnection voltage.
                                 Carries its own differential protection
                                 zone. \\
    8  & Generator breaker 52G, POI
                               & Switching point between island and
                                 grid-connected operation, closed only on
                                 synchroniser permission. Revenue metering and
                                 the regulated POI quantities sit here. \\
    9  & Air-cooled condenser  & Rejects cycle heat to ambient; VFD-driven fan
                                 motors set condensing pressure. A cooling-tower
                                 loop replaces it on water-cooled plants. \\
    10 & Working-fluid feed pump
                               & VFD-driven pump closing the loop. Its speed is
                                 the principal handle on working-fluid mass
                                 flow, and therefore on plant output over
                                 minutes. \\
    \bottomrule
  \end{tabularx}
  \caption{Physical parts of the ORC energy path, keyed to
           Figure~\ref{fig:orc-control-cycle}.}
  \label{tab:orc-parts}
\end{table}

\subsubsection{Control loops}
Each loop pairs a sensor, a controller and an actuator around one of the parts
above. The fast loops (milliseconds to seconds) reside in dedicated hardware ---
governor, AVR, protection --- while the slow thermal loops (seconds to minutes)
reside in the plant DCS. Only the fast loops survive into the dynamic model of
Section~\ref{sec:model}.

\begin{table}[htbp]
  \centering
  \small
  \renewcommand{\arraystretch}{1.25}
  \begin{tabularx}{\linewidth}{@{}p{2.1cm} p{3.2cm} p{3.7cm} X@{}}
    \toprule
    \textbf{Loop} & \textbf{Measured} & \textbf{Actuator} &
    \textbf{Regulated quantity} \\
    \midrule
    Speed / load         & Shaft speed, MW output, turbine inlet pressure
                         & TCV (3)
                         & Speed during run-up; MW or inlet pressure once
                           synchronised \\
    Evaporator pressure  & Evaporator pressure and superheat (2)
                         & Feed-pump VFD (10)
                         & Vapour conditions at turbine inlet as heat input
                           varies \\
    Condensing pressure  & Condenser pressure (9)
                         & ACC fan VFDs (9)
                         & Back-pressure near optimum across ambient swings \\
    Voltage / VAR        & Generator terminal voltage and current (6)
                         & Exciter field current (6)
                         & Terminal voltage, or the power-factor / MVAr
                           setpoint \\
    Synchronisation      & Voltage, frequency and phase across 52G
                         & Governor and AVR bias, 52G close coil (8)
                         & Match across the breaker before closing is
                           permitted \\
    Plant / POI          & MW, MVAr and voltage at the POI meter
                         & Setpoints to all controllers above
                         & Dispatch target and the POI voltage / power-factor
                           schedule \\
    \bottomrule
  \end{tabularx}
  \caption{Control loops of the ORC power block. Numbers in parentheses refer
           to Table~\ref{tab:orc-parts}.}
  \label{tab:orc-loops}
\end{table}

\subsection{Protection}
Protection acts once and latches rather than regulating continuously. The
generator protection relay --- differential (87G), reverse power (32), loss of
field (40), over/under-frequency (81) and over/under-voltage (27/59) --- trips
the 52G breaker and simultaneously commands the turbine trip system, which dumps
hydraulic pressure and allows the spring-closed ESV to shut. The reverse-power
element is of particular relevance to an ORC unit: left connected without
working-fluid flow, the machine would motor the generator and drive the
expander backwards.

An independent electronic overspeed system trips the ESV directly, without
passing through the DCS, so that loss of the control system can never disable
overspeed protection. The GSU transformer carries its own differential and
overcurrent zone.

In the simulation model these functions appear as breaker events and relay
settings, not as dynamics.

\subsection{Positive-sequence dynamic model}
\label{sec:model}
In a positive-sequence dynamic model (PSS/E, PSLF, PowerFactory, ETAP) the whole
plant of Figure~\ref{fig:orc-control-cycle} collapses to three interconnected
blocks plus the network solution. The thermal loops --- evaporator pressure,
condensing pressure, feed pump and fans --- are far too slow to influence
stability timeframes and disappear into the boundary condition
$P_{\mathrm{mech}} \approx \mathrm{const}$. What remains is the wiring of
Figure~\ref{fig:model-wiring}, which is identical in every tool; only the block
names change.

\begin{figure}[htbp]
  \centering
  \resizebox{0.96\linewidth}{!}{%
  \begin{tikzpicture}[
    font=\scriptsize,
    blk/.style={draw=atFg, line width=0.6pt, fill=atBgSubtle, rounded corners=1pt,
                align=center, inner sep=3pt},
    sblk/.style={draw=atAccentFg, line width=0.7pt, dashed, fill=white,
                 rounded corners=1pt, align=center, inner sep=3pt},
    sumj/.style={draw=atFg, line width=0.6pt, fill=white, circle, inner sep=1.5pt},
    sig/.style={draw=atFg, line width=0.9pt, -{Latex[length=2mm]}},
    fb/.style={draw=atAccentFg, dashed, line width=0.6pt, -{Latex[length=1.8mm]}},
    lab/.style={font=\tiny, text=atFgMuted},
    sg/.style={font=\tiny, text=atAccentFg},
    tap/.style={circle, fill=atAccentFg, inner sep=1.1pt},
  ]
  \node[blk, minimum width=28mm, minimum height=11mm] (tgov) at (0,1.4)
        {\texttt{TGOV1}\\[1pt]\tiny turbine--governor};
  \node[blk, minimum width=28mm, minimum height=11mm] (sexs) at (0,-0.8)
        {\texttt{SEXS} $\rightarrow$ \texttt{AC7B}\\[1pt]\tiny AVR + exciter};
  \node[blk, minimum width=26mm, minimum height=34mm] (gen) at (4.6,0.3)
        {\texttt{GENROU}\\[2pt]\tiny synchronous\\\tiny machine\\[2pt]
         \tiny(\texttt{GENSAL} if\\\tiny salient-pole)};
  \node[blk, minimum width=28mm, minimum height=18mm] (net) at (8.2,0.3)
        {NETWORK\\[1pt]\tiny GSU impedance + $Y_{\mathrm{bus}}$\\\tiny POI / grid};
  \node[sumj] (sum) at (-2.8,-0.8) {$\Sigma$};
  \node[sblk, minimum width=30mm, minimum height=10mm] (pss) at (-0.6,-2.6)
        {\texttt{PSS2B} (WECC: required)\\[1pt]\tiny input $\Delta\omega$,
         output $V_s$};

  \draw[sig] (-5.0,1.4) -- (tgov.west);
  \node[lab, anchor=south] at (-4.2,1.5) {$P_{\mathrm{ref}}$};
  \draw[sig] (-5.0,-0.8) -- (sum.west);
  \node[lab, anchor=south] at (-4.2,-0.7) {$V_{\mathrm{ref}}$};
  \node[lab] at (-3.1,-0.42) {$+$};
  \node[lab] at (-3.1,-1.20) {$-$};
  \draw[sig] (sum.east) -- (sexs.west);
  \draw[sig] (tgov.east) -- (3.3,1.4);
  \node[lab, anchor=south] at (2.35,1.5) {$P_{\mathrm{mech}}$};
  \draw[sig] (sexs.east) -- (3.3,-0.8);
  \node[lab, anchor=south] at (2.35,-0.7) {$E_{\mathrm{fd}}$};
  \draw[draw=atFg, line width=0.9pt, {Latex[length=2mm]}-{Latex[length=2mm]}]
        (gen.east) -- (net.west);
  \node[lab, anchor=south] at (6.35,0.42) {$V_t$, $I_t$};

  % feedbacks
  \draw[fb] (4.6,2.0) -- (4.6,2.9) -- (0.7,2.9) -- (0.7,1.95);
  \node[sg, anchor=south] at (2.65,2.95) {$\Delta\omega$ (rotor speed deviation)};
  \node[tap] at (6.35,0.3) {};
  \draw[fb] (6.35,0.18) -- (6.35,-3.5) -- (-2.8,-3.5) -- (-2.8,-1.08);
  \node[sg, anchor=south] at (2.4,-3.45) {$V_t$ (terminal voltage)};
  \draw[fb] (-2.1,-2.6) -- (-3.7,-2.6) -- (-3.7,-0.2) -- (-2.8,-0.2) -- (-2.8,-0.52);
  \node[sg, anchor=west] at (-3.60,-0.05) {$V_s$};
  \node[lab, anchor=north, align=center] at (0,0.72)
        {base-loaded: $P_{\mathrm{ref}}$ fixed,\\or unit flagged non-responsive};
  \end{tikzpicture}}
  \caption{Model wiring implemented by every positive-sequence tool. Solid lines
  are model signals; dashed lines are the two feedbacks that close the loops ---
  $\Delta\omega$ to the governor and $V_t$ to the AVR summing junction.
  Protection and the synchroniser do not appear as dynamics.}
  \label{fig:model-wiring}
\end{figure}

\subsection{Turbine--governor model}
The governor block of Figure~\ref{fig:model-wiring} is expanded in
Figure~\ref{fig:tgov1}. \texttt{TGOV1} is the model named by the NERC
turbine-governor modelling guideline for a simple, non-reheat thermal unit and
is the appropriate choice for an ORC turbine.

\begin{figure}[htbp]
  \centering
  \resizebox{0.94\linewidth}{!}{%
  \begin{tikzpicture}[
    font=\scriptsize,
    blk/.style={draw=atFg, line width=0.6pt, fill=atBgSubtle, rounded corners=1pt,
                align=center, inner sep=3pt},
    sumj/.style={draw=atFg, line width=0.6pt, fill=white, circle, inner sep=1.5pt},
    sig/.style={draw=atFg, line width=0.9pt, -{Latex[length=2mm]}},
    sigp/.style={draw=atFg, line width=0.9pt},
    lab/.style={font=\tiny, text=atFgMuted},
    tap/.style={circle, fill=atFg, inner sep=1.1pt},
  ]
  \node[sumj] (s1) at (-2.2,0) {$\Sigma$};
  \node[blk, minimum width=22mm, minimum height=12mm] (lag) at (0.4,0)
        {$\dfrac{1}{1+sT_1}$};
  \node[blk, minimum width=24mm, minimum height=12mm] (ll) at (3.4,0)
        {$\dfrac{1+sT_2}{1+sT_3}$};
  \node[sumj] (s2) at (5.9,0) {$\Sigma$};
  \node[blk, minimum width=12mm, minimum height=9mm] (rr) at (-3.4,-2.0) {$1/R$};
  \node[blk, minimum width=12mm, minimum height=9mm] (dt) at (4.3,-3.0) {$D_t$};

  \draw[sig] (-4.4,0) -- (s1.west);
  \node[lab, anchor=south] at (-3.7,0.10) {$P_{\mathrm{ref}}$};
  \node[lab] at (-2.5,0.30) {$+$};
  \draw[sig] (s1.east) -- (lag.west);
  \draw[sig] (lag.east) -- (ll.west);
  \draw[sig] (ll.east) -- (s2.west);
  \draw[sig] (s2.east) -- (8.0,0);
  \node[lab, anchor=south] at (7.2,0.10) {$P_{\mathrm{mech}}$};
  \node[lab, anchor=south] at (0.4,0.66) {$V_{\max}$};
  \node[lab, anchor=north] at (0.4,-0.66) {$V_{\min}$};

  \draw[sigp] (-4.9,-2.0) -- (rr.west);
  \node[lab, anchor=south] at (-4.55,-1.90) {$\Delta\omega$};
  \node[tap] (j) at (-4.3,-2.0) {};
  \draw[sig] (rr.east) -- (-2.2,-2.0) -- (-2.2,-0.28);
  \node[lab] at (-1.95,-0.55) {$-$};
  \draw[sigp] (-4.3,-2.0) -- (-4.3,-3.0);
  \draw[sig] (-4.3,-3.0) -- (dt.west);
  \draw[sig] (dt.east) -- (5.9,-3.0) -- (5.9,-0.28);
  \node[lab] at (6.15,-0.55) {$-$};
  \end{tikzpicture}}
  \caption{\texttt{TGOV1} turbine--governor model. For an ORC unit $T_1$ lumps
  the valve stroke and the evaporator vapour-volume lag ($\approx$
  0.4--1.5\,s); there is no reheater, so $T_2 = T_3$ and the lead--lag drops
  out. Base-loaded operation is represented either by
  $V_{\max} = V_{\min} = P_0$ or by flagging the unit non-responsive.}
  \label{fig:tgov1}
\end{figure}

\subsection{Excitation system model}
The excitation block of Figure~\ref{fig:model-wiring} is expanded in
Figure~\ref{fig:sexs}, shown here in its screening-level \texttt{SEXS} form.

\begin{figure}[htbp]
  \centering
  \resizebox{0.84\linewidth}{!}{%
  \begin{tikzpicture}[
    font=\scriptsize,
    blk/.style={draw=atFg, line width=0.6pt, fill=atBgSubtle, rounded corners=1pt,
                align=center, inner sep=3pt},
    sumj/.style={draw=atFg, line width=0.6pt, fill=white, circle, inner sep=1.5pt},
    sig/.style={draw=atFg, line width=0.9pt, -{Latex[length=2mm]}},
    lab/.style={font=\tiny, text=atFgMuted},
  ]
  \node[sumj] (s) at (-2.0,0) {$\Sigma$};
  \node[blk, minimum width=24mm, minimum height=12mm] (ll) at (0.7,0)
        {$\dfrac{1+sT_A}{1+sT_B}$};
  \node[blk, minimum width=22mm, minimum height=12mm] (k) at (3.7,0)
        {$\dfrac{K}{1+sT_E}$};
  \draw[sig] (-4.2,0) -- (s.west);
  \node[lab, anchor=south] at (-3.5,0.10) {$V_{\mathrm{ref}}$};
  \draw[sig] (-2.0,1.7) -- (s.north);
  \node[lab, anchor=west] at (-1.90,1.35) {$V_s$ (PSS) \ $+$};
  \draw[sig] (-2.0,-1.7) -- (s.south);
  \node[lab, anchor=west] at (-1.90,-1.35) {$V_t$ \ $-$};
  \node[lab] at (-2.30,0.30) {$+$};
  \draw[sig] (s.east) -- (ll.west);
  \draw[sig] (ll.east) -- (k.west);
  \draw[sig] (k.east) -- (6.6,0);
  \node[lab, anchor=south] at (5.8,0.10) {$E_{\mathrm{fd}}$};
  \node[lab, anchor=south] at (3.7,0.66) {$E_{\max}$};
  \node[lab, anchor=north] at (3.7,-0.66) {$E_{\min}$};
  \end{tikzpicture}}
  \caption{\texttt{SEXS} excitation model: voltage error, regulator lead--lag,
  exciter gain with ceiling limits. It is a screening-level placeholder; once
  OEM excitation data is available the block is replaced by IEEE
  \texttt{AC7B}/\texttt{AC8B} for a brushless exciter (the usual case on ORC
  machines) or \texttt{ST4B} for a static exciter. The wiring of
  Figure~\ref{fig:model-wiring} is unchanged.}
  \label{fig:sexs}
\end{figure}

\subsection{Power system stabilizer and WECC requirement}
\label{sec:pss}
A power system stabilizer (PSS) is a supplementary controller added to the
excitation system --- the dashed block in Figure~\ref{fig:model-wiring}. The AVR
alone regulates terminal voltage, but a fast, high-gain AVR erodes the damping
of the rotor's electromechanical oscillations. Following a disturbance every
synchronous rotor swings against the rest of the system at 0.1--3\,Hz: local
machine-against-system modes around 1--2\,Hz, and inter-area modes, in which
large groups of machines in one region swing against another, at 0.1--0.8\,Hz.
The PSS measures rotor speed deviation and electrical power --- the modern
dual-input integral-of-accelerating-power design, models \texttt{PSS2B} and
\texttt{PSS2C} --- phase-compensates that signal and injects it as $V_s$ into
the AVR summing junction, so that the field is modulated in phase with rotor
speed and produces a damping torque. The PSS adds damping; it does not alter
steady-state voltage control.

The Western Interconnection has long transmission corridors and well-known
inter-area modes --- the $\sim$0.25\,Hz North--South mode was central to the
August 1996 west-wide separation --- and WECC is consequently the one region
with its own enforceable regional PSS standard.

\paragraph{VAR-501-WECC-4 requirements.} The standard applies to synchronous
generators connected to the Bulk Electric System that meet the definition of
Commercial Operation; a 40\,MW / 47\,MVA unit falls within scope. The
obligations relevant to this project are:

\begin{itemize}\setlength{\itemsep}{0.25em}
  \item \textbf{Installation} --- PSS installed and placed in service within
        180 days of a new generator connecting to the BES, or of the
        replacement of an existing voltage regulator.
  \item \textbf{In service} --- the PSS must be in service whenever the unit is
        synchronised, other than during component failure, testing or
        maintenance, or as otherwise agreed with the Transmission Operator.
        Outages shorter than 30 minutes are not violations; inoperable
        equipment must be repaired or replaced within 24 months.
  \item \textbf{Tuning} --- phase response within $\pm$30$^\circ$ over the
        defined frequency range; output authority of at least $\pm$5\,\% of
        terminal voltage; gain set between one-third and one-half of the
        maximum practical gain; washout time constant not exceeding 30\,s.
  \item \textbf{Exemption} --- narrow, and limited to excitation systems
        incapable of meeting the tuning requirements; it lapses when the
        voltage regulator is replaced or retrofitted.
\end{itemize}

\noindent
The practical consequences for this project are that the PSS must be specified
together with the excitation system at the time of purchase, that commissioning
must include PSS tuning tests, and that a \texttt{PSS2B}/\texttt{PSS2C} record
must accompany the exciter model in the submitted dynamics data, populated from
those tests.

\subsection{Model parameters}
Table~\ref{tab:dyn-params} gives screening-level values on a 47\,MVA machine
base (40\,MW at 0.85 power factor), 60\,Hz, four-pole through gearbox. The
characteristically ORC parameter is the \textbf{low inertia}: a compact expander
and gearbox store far less rotational energy than a steam turbine train, so $H$
sits at the bottom of the synchronous range.

\begin{table}[htbp]
  \centering
  \small
  \renewcommand{\arraystretch}{1.25}
  \begin{tabularx}{\linewidth}{@{}l X@{}}
    \toprule
    \textbf{Model} & \textbf{Typical parameters} \\
    \midrule
    \texttt{GENROU}
      & $T'_{d0}=5.0$\,s, $T''_{d0}=0.05$, $T'_{q0}=0.8$, $T''_{q0}=0.08$,
        $H=1.8$\,s, $D=0$, $X_d=1.80$, $X_q=1.70$, $X'_d=0.30$, $X'_q=0.55$,
        $X''=0.20$, $X_l=0.12$, $S(1.0)=0.10$, $S(1.2)=0.40$. \newline
        $H \approx 1.5$--$2.5$\,s is typical for ORC turbogenerators. Use
        \texttt{GENSAL}/\texttt{GENTPJ} if the OEM machine is salient-pole,
        which is common at 1{,}500/1{,}800\,rpm. \\
    \texttt{SEXS}
      & $T_A/T_B=0.10$, $T_B=10.0$\,s, $K=100$, $T_E=0.10$\,s,
        $E_{\min}=0.0$, $E_{\max}=4.0$. Placeholder until the OEM AVR data
        sheet supports \texttt{AC7B} or \texttt{ST4B}. \\
    \texttt{TGOV1}
      & $R=0.05$, $T_1=0.40$\,s, $V_{\max}=1.0$, $V_{\min}=0.0$, $T_2=1.0$,
        $T_3=1.0$, $D_t=0$. For a base-loaded unit set
        $V_{\max}=V_{\min}=P_0$, or omit the governor and flag the unit
        non-responsive. \\
    \bottomrule
  \end{tabularx}
  \caption{Screening-level dynamic model parameters, 40\,MW ORC unit on a
           47\,MVA base. Replace with OEM data before any interconnection
           filing.}
  \label{tab:dyn-params}
\end{table}

\noindent
The corresponding PSS/E dynamics record is:

{\small
\begin{verbatim}
/ 40 MW ORC unit at bus 101, id 1 - machine base 47.0 MVA, 60 Hz
/ Placeholder screening data: replace with OEM values for any filing
101 'GENROU' 1   5.00  0.05  0.80  0.08  1.80  0.00
                 1.80  1.70  0.30  0.55  0.20  0.12  0.10  0.40 /
101 'SEXS'   1   0.10  10.0  100.0 0.10  0.00  4.00 /
101 'TGOV1'  1   0.05  0.40  1.00  0.00  1.00  1.00  0.00 /
\end{verbatim}}

\noindent
A \texttt{PSS2B} record must be added alongside the exciter before submission,
per Section~\ref{sec:pss}.

\subsection{Portability across tools}
The same three blocks exist under the same or near-identical names in every
environment:

\begin{itemize}\setlength{\itemsep}{0.25em}
  \item \textbf{PSLF} --- \texttt{genrou} / \texttt{sexs} / \texttt{tgov1},
        with an \texttt{epcl} equivalent of the \texttt{.dyr} record.
  \item \textbf{PowerFactory, ETAP} --- shipped as standard IEEE model
        implementations; parameters transfer directly from
        Table~\ref{tab:dyn-params}.
  \item \textbf{MATLAB/Simulink} --- Figures~\ref{fig:tgov1} and
        \ref{fig:sexs} are wired directly from transfer-function blocks around
        the built-in synchronous machine model.
  \item \textbf{PSCAD (EMT)} --- retain the machine and exciter and drive the
        shaft with constant mechanical torque. The ORC thermal side contributes
        nothing at EMT timescales.
\end{itemize}

\subsection{Open items}
\begin{itemize}\setlength{\itemsep}{0.25em}
  \item Obtain the OEM generator data sheet and confirm round-rotor versus
        salient-pole construction, which selects \texttt{GENROU} against
        \texttt{GENSAL}/\texttt{GENTPJ}.
  \item Obtain the excitation system type and data to replace \texttt{SEXS}
        with \texttt{AC7B}/\texttt{AC8B} or \texttt{ST4B}.
  \item Confirm whether the unit will provide primary frequency response; if
        not, agree the non-responsive representation with the Transmission
        Planner.
  \item Specify the PSS with the excitation system and schedule commissioning
        tuning tests (Section~\ref{sec:pss}).
  \item Obtain GSU impedance, machine saturation data and the reactive
        capability curve. These bind the model to the real plant more tightly
        than any control parameter.
\end{itemize}
